% kotexcweb.tex -- Korean localization and a LuaTeX PDF back end for CWEB.
%
% Put  \input kotexcweb  as the FIRST line of the limbo of a .w file and process
% the woven output with luatex (not pdftex):  luatex foo.tex.  cweave emits
% "\input cwebmac" as the very first line of foo.tex, so this file is read
% afterwards and overrides cwebmac wherever Korean text or the LuaTeX engine
% needs it.  It has to come first in limbo because it reads cwebmac a second
% time (see the next paragraph), which would undo anything defined in limbo
% before it -- \datethis, \def\title, \pageheight and the like all belong after
% the \input, not before it.
%
% Requires luatexko and the Noto Serif/Sans CJK KR fonts.  The hangul font setup
% is self-contained below; edit the \sethangulfont lines to change typefaces.
%
% What cweave does to Korean before TeX ever sees it.  It reads a line of a .w
% file into a 100-byte buffer, and a Korean character is three bytes, so a
% source line holds barely thirty characters before cweave stops with "Input
% line too long" -- and a C string on a line that overruns the buffer comes out
% garbled, half a character landing in the index.  The change files next to
% this one, comm-ko.ch with ctang-ko.ch and cweav-ko.ch, raise the buffer to a
% kilobyte and both symptoms go away:
%
%   make CCHANGES=comm-ko.ch TCHANGES=ctang-ko.ch WCHANGES=cweav-ko.ch
%
% Without them, keep Korean source lines short.  (cweave 4 reads longer lines
% to begin with, and this file digests its output too, as long as the cwebmac
% it finds is still the 3.70 one.)  What cweave does inside a string in any
% case -- a discretionary break every twenty bytes, counted in bytes, which
% cuts Korean characters in half -- is repaired below.

\ifx\pdfextension\undefined
  \errmessage{kotexcweb.tex needs LuaTeX; run luatex, not pdftex or tex}\fi
% This file follows the cwebmac.tex that comes with CWEB 3 (version 3.70, the
% one next to it in this directory).  CWEB 4's cwebmac is a different macro
% file: it knows about LuaTeX already, and it renamed \ifacro to \ifacrohint,
% so the overrides below would not fit it.  TeX looks in the current directory
% first, so weave and typeset in the directory that holds this cwebmac.
\ifx\cwebversion\undefined\else\ifnum\cwebversion>3
  \errmessage{kotexcweb.tex expects cwebmac 3.x, but cwebmac
    \number\cwebversion.\number\cwebrevision\space was loaded}\fi\fi

% ---- pdfTeX's primitives, and a second reading of cwebmac -------------------
% cwebmac decides once and for all, at the moment it is read, whether to emit
% hyperlinks and bookmarks: it looks for \pdfoutput, and LuaTeX has no such
% primitive (it became \outputmode in LuaTeX 0.85).  So \ifacro comes out false
% and cwebmac skips its entire block of PDF macros -- \pdflink, \pdfnote,
% \pdfURL and the ``sanitizer'' that turns a section name into a bookmark title.
% Nothing can switch that block on afterwards, and we cannot get in ahead of
% cwebmac either, since cweave writes the \input line itself.
%
% What we can do is make LuaTeX look like pdfTeX and read cwebmac again.  The
% six primitives cwebmac uses are all spellings of LuaTeX's \pdfextension, and
% \pdfoutput is just a counter to cwebmac (it tests it and sets it to 1).
%
% The second reading repeats cwebmac's assignments and allocates a few more
% registers, which costs nothing.  Its last line, though, contributes a
% \vbox to\vsize{} -- there so that the first \topmark will not be null -- and a
% second one of those would be shipped out as an empty first page, since the
% ``the first page is garbage'' output routine throws away one page, not two.
% So \vbox is booby-trapped just long enough to catch that box in a register:
% it is the only \vbox cwebmac executes while it is being read, and it puts
% itself back as soon as it has caught it.
\newcount\pdfoutput \pdfoutput=1
\protected\def\pdfliteral{\pdfextension literal }
\protected\def\pdfdest{\pdfextension dest }
\protected\def\pdfoutline{\pdfextension outline }
\protected\def\pdfcatalog{\pdfextension catalog }
\protected\def\pdfannotlink{\pdfextension startlink }
\protected\def\pdfendlink{\pdfextension endlink\relax} % \relax: no stray space
% luatexko comes first, because it loads atbegshi and hence iftex, and iftex
% declares an \ifpdftex of its own -- false under LuaTeX, quite correctly, but
% it would shadow cwebmac's flag of the same name and switch the links and
% bookmarks back off.  Reading cwebmac afterwards gives cwebmac the last word.
\input luatexko.sty
\newbox\spareframe
\let\korealvbox=\vbox \def\vbox{\let\vbox=\korealvbox\setbox\spareframe=\vbox}
\input cwebmac

% ---- Korean fonts -----------------------------------------------------------
% luatexko provides Hangul shaping (and automatic fallback, so Korean appears
% even through cwebmac's CM fonts).  The plain-TeX hangul setup is borrowed from
% luatexko-doc.tex: a serif body (Noto Serif), a sans running head (Noto Sans),
% and a monospace face (Noto Sans Mono) paired with the matching CM fonts.
%
% Sizes.  Hangul fills nearly its whole em box while CM's lowercase reaches only
% an unusually small x-height (0.431em, against 0.511em for the Latin that Noto
% itself ships), so Hangul set at the nominal size looks too big next to CM.
% Matching the ink each script actually spans -- CM's ascender-to-descender is
% 0.889em, an average Hangul syllable 0.924em -- gives 0.962 of the Latin size;
% against cmtt, whose span is only 0.833em, the monospace face wants 0.93.  So
% each Hangul font below is (its CM partner's design size) times that ratio.  CM
% keeps the same proportions at every design size (x-height is 0.431em in cmr10,
% cmr9 and cmr8 alike), which is why a constant ratio is right and the fixed
% -0.2pt offset used before was not: it drifted from .980 at ten point to .971
% at seven.  The interhangul/interlatincjk/charraise settings are all in em and
% therefore follow the size on their own; they are not restated per size.

\sethangulfont\tenmj="Noto Serif CJK KR Light:%
  language=KOR;%
  script=hang;%
  interhangul=-0.04em;%
  interlatincjk=.125em;%
  charraise=-0.06em;%
  +compresspunctuations;%
  expansion=default;" at 9.62pt
\sethangulfont\tentz="Noto Sans Mono CJK KR:%
  language=KOR;%
  script=hang;%
  interhangul=-0.04em;%
  interlatincjk=.125em;%
  charraise=-0.03em;%
  +compresspunctuations;%
  expansion=default;" at 9.3pt
\sethangulfont\tenbd="Noto Serif CJK KR Bold:%
  language=KOR;%
  script=hang;%
  interhangul=-0.04em;%
  interlatincjk=.125em;%
  charraise=-0.06em;%
  +compresspunctuations;%
  expansion=default;" at 9.62pt
\sethangulfont\tengt="Noto Sans CJK KR Light:%
  language=KOR;%
  script=hang;%
  interhangul=-0.04em;%
  interlatincjk=.125em;%
  charraise=-0.06em;%
  +compresspunctuations;%
  expansion=default;" at 9.62pt

\sethangulfont\ninemj="Noto Serif CJK KR Light:%
  language=KOR;%
  script=hang;%
  interhangul=-0.04em;%
  interlatincjk=.125em;%
  charraise=-0.06em;%
  +compresspunctuations;%
  expansion=default;" at 8.66pt
\sethangulfont\ninetz="Noto Sans Mono CJK KR:%
  language=KOR;%
  script=hang;%
  interhangul=-0.04em;%
  interlatincjk=.125em;%
  charraise=-0.03em;%
  +compresspunctuations;%
  expansion=default;" at 8.37pt
\sethangulfont\ninebd="Noto Serif CJK KR Bold:%
  language=KOR;%
  script=hang;%
  interhangul=-0.04em;%
  interlatincjk=.125em;%
  charraise=-0.06em;%
  +compresspunctuations;%
  expansion=default;" at 8.66pt
\sethangulfont\ninegt="Noto Sans CJK KR Light:%
  language=KOR;%
  script=hang;%
  interhangul=-0.04em;%
  interlatincjk=.125em;%
  charraise=-0.06em;%
  +compresspunctuations;%
  expansion=default;" at 8.66pt

\sethangulfont\eightmj="Noto Serif CJK KR Light:%
  language=KOR;%
  script=hang;%
  interhangul=-0.04em;%
  interlatincjk=.125em;%
  charraise=-0.06em;%
  +compresspunctuations;%
  expansion=default;" at 7.7pt
\sethangulfont\eighttz="Noto Sans Mono CJK KR:%
  language=KOR;%
  script=hang;%
  interhangul=-0.04em;%
  interlatincjk=.125em;%
  charraise=-0.03em;%
  +compresspunctuations;%
  expansion=default;" at 7.44pt
\sethangulfont\eightbd="Noto Serif CJK KR Bold:%
  language=KOR;%
  script=hang;%
  interhangul=-0.04em;%
  interlatincjk=.125em;%
  charraise=-0.06em;%
  +compresspunctuations;%
  expansion=default;" at 7.7pt
\sethangulfont\eightgt="Noto Sans CJK KR Light:%
  language=KOR;%
  script=hang;%
  interhangul=-0.04em;%
  interlatincjk=.125em;%
  charraise=-0.06em;%
  +compresspunctuations;%
  expansion=default;" at 7.7pt

\sethangulfont\sevenmj="Noto Serif CJK KR Light:%
  language=KOR;%
  script=hang;%
  interhangul=-0.04em;%
  interlatincjk=.125em;%
  charraise=-0.06em;%
  +compresspunctuations;%
  expansion=default;" at 6.73pt

% hangul in math mode (identifiers and section names inside |...| are set in
% math mode, so a Korean section name needs these).
\font\texthangul="Noto Serif CJK KR Regular" at 9.62pt
\font\texthangulnine="Noto Serif CJK KR Regular" at 8.66pt
\font\texthanguleight="Noto Serif CJK KR Regular" at 7.7pt
\font\scripthangul="Noto Serif CJK KR Regular" at 6.73pt
\font\scripthangulsix="Noto Serif CJK KR Regular" at 5.77pt
\font\scriptscripthangul="Noto Serif CJK KR Medium" at 4.81pt

% ---- size switches ----------------------------------------------------------
% Size-switch macros (borrowed from manmac.tex) that bind each CM family to its
% matching hangul font, so \rm/\it/\bf/\tt carry Korean as well as Latin.
% cwebmac declares \ninerm and \eightrm as fonts and uses them as font switches
% (in \mc, \sc, \note, \X, the running heads and the index); the fonts are kept
% here under the names \ninermfont and \eightrmfont, and the switches are
% redefined below to pull in the matching hangul.
\catcode`@=11 % borrow the private macros of PLAIN (with care)

\font\ninermfont=cmr9
\font\eightrmfont=cmr8
\font\sixrm=cmr6

\font\ninei=cmmi9
\font\eighti=cmmi8
\font\sixi=cmmi6
\skewchar\ninei='177 \skewchar\eighti='177 \skewchar\sixi='177

\font\ninesy=cmsy9
\font\eightsy=cmsy8
\font\sixsy=cmsy6
\skewchar\ninesy='60 \skewchar\eightsy='60 \skewchar\sixsy='60

\font\ninebf=cmbx9
\font\eightbf=cmbx8
\font\sixbf=cmbx6

\font\ninett=cmtt9
\font\eighttt=cmtt8
\hyphenchar\tentt=-1 % inhibit hyphenation in typewriter type
\hyphenchar\ninett=-1
\hyphenchar\eighttt=-1

\font\ninesl=cmsl9
\font\eightsl=cmsl8

\font\nineit=cmti9
\font\eightit=cmti8

\newskip\ttglue
\def\tenpoint{\def\rm{\fam0\tenrm\tenmj}%
  \textfont0=\tenrm \scriptfont0=\sevenrm \scriptscriptfont0=\fiverm
  \textfont1=\teni \scriptfont1=\seveni \scriptscriptfont1=\fivei
  \textfont2=\tensy \scriptfont2=\sevensy \scriptscriptfont2=\fivesy
  \textfont3=\tenex \scriptfont3=\tenex \scriptscriptfont3=\tenex
  \setmathhangulfonts\texthangul\scripthangul\scriptscripthangul
  \def\it{\fam\itfam\tenit\tengt}%
  \textfont\itfam=\tenit
  \def\sl{\fam\slfam\tensl\tengt}%
  \textfont\slfam=\tensl
  \def\bf{\fam\bffam\tenbf\tenbd}%
  \textfont\bffam=\tenbf \scriptfont\bffam=\sevenbf
  \scriptscriptfont\bffam=\fivebf
  \def\tt{\fam\ttfam\tentt\tentz}%
  \textfont\ttfam=\tentt
  \tt \ttglue=.5em plus.25em minus.15em
  \normalbaselineskip=14pt
  \let\big=\tenbig
  \setbox\strutbox=\hbox{\vrule height8.5pt depth3.5pt width\z@}%
  \normalbaselines\rm}

\def\ninepoint{\def\rm{\fam0\ninermfont\ninemj}%
  \textfont0=\ninermfont \scriptfont0=\sixrm \scriptscriptfont0=\fiverm
  \textfont1=\ninei \scriptfont1=\sixi \scriptscriptfont1=\fivei
  \textfont2=\ninesy \scriptfont2=\sixsy \scriptscriptfont2=\fivesy
  \textfont3=\tenex \scriptfont3=\tenex \scriptscriptfont3=\tenex
  \setmathhangulfonts\texthangulnine\scripthangulsix\scriptscripthangul
  \def\it{\fam\itfam\nineit\ninegt}%
  \textfont\itfam=\nineit
  \def\sl{\fam\slfam\ninesl\ninegt}%
  \textfont\slfam=\ninesl
  \def\bf{\fam\bffam\ninebf\ninebd}%
  \textfont\bffam=\ninebf \scriptfont\bffam=\sixbf
   \scriptscriptfont\bffam=\fivebf
  \def\tt{\fam\ttfam\ninett\ninetz}%
  \textfont\ttfam=\ninett
  \tt \ttglue=.5em plus.25em minus.15em
  \normalbaselineskip=13pt
  \let\big=\ninebig
  \setbox\strutbox=\hbox{\vrule height8pt depth3pt width\z@}%
  \normalbaselines\rm}

\def\eightpoint{\def\rm{\fam0\eightrmfont\eightmj}%
  \textfont0=\eightrmfont \scriptfont0=\sixrm \scriptscriptfont0=\fiverm
  \textfont1=\eighti \scriptfont1=\sixi \scriptscriptfont1=\fivei
  \textfont2=\eightsy \scriptfont2=\sixsy \scriptscriptfont2=\fivesy
  \textfont3=\tenex \scriptfont3=\tenex \scriptscriptfont3=\tenex
  \setmathhangulfonts\texthanguleight\scripthangulsix\scriptscripthangul
  \def\it{\fam\itfam\eightit\eightgt}%
  \textfont\itfam=\eightit
  \def\sl{\fam\slfam\eightsl\eightgt}%
  \textfont\slfam=\eightsl
  \def\bf{\fam\bffam\eightbf\eightbd}%
  \textfont\bffam=\eightbf \scriptfont\bffam=\sixbf
   \scriptscriptfont\bffam=\fivebf
  \def\tt{\fam\ttfam\eighttt\eighttz}%
  \textfont\ttfam=\eighttt
  \tt \ttglue=.5em plus.25em minus.15em
  \normalbaselineskip=11pt
  \let\big=\eightbig
  \setbox\strutbox=\hbox{\vrule height7pt depth2pt width\z@}%
  \normalbaselines\rm}

\def\tenmath{\tenpoint\fam-1 } % use after $ in ninepoint sections
\def\tenbig#1{{\hbox{$\left#1\vbox to8.5pt{}\right.\n@space$}}}
\def\ninebig#1{{\hbox{$\textfont0=\tenrm\textfont2=\tensy
  \left#1\vbox to7.25pt{}\right.\n@space$}}}
\def\eightbig#1{{\hbox{$\textfont0=\ninermfont\textfont2=\ninesy
  \left#1\vbox to6.5pt{}\right.\n@space$}}}

\catcode`\@=12
\tenpoint

% ---- cwebmac's font hooks ---------------------------------------------------
% Each of these was \let to a bare CM font by cwebmac, which leaves Korean to
% luatexko's fallback -- i.e. at whatever size and shape the fallback happens to
% pick.  Make them switch the hangul companion too.
\def\eightrm{\eightrmfont\eightmj} % notes, index, headers: cmr8 + 7.8pt serif
\def\ninerm{\ninermfont\ninemj}
\def\mc{\ninerm}                   % medium caps
\def\sc{\eightrm}                  % smallish caps
\def\mainfont{\tenrm\tenmj}        % page numbers and section numbers in headers
\def\cmntfont{\tenrm\tenmj}        % comments in the C parts
\def\smallfont{\eightrm}           % the small print, as a name of its own

% The title fonts: keep cwebmac's Latin faces for Latin and punctuation and pair
% a Hangul companion of the same size, so only Korean characters take the CJK
% face.  cwebmac has \titlefont=cmr7 scaled\magstep4 (~14.5pt) for the contents
% page and \ttitlefont=cmtt10 scaled\magstep2 (~14.4pt) for typewriter text in a
% title; the companions take the same ratios the body fonts do, 0.962 of the
% serif partner and 0.93 of the typewriter one.
\let\latintitlefont=\titlefont
\sethangulfont\hangultitlefont="Noto Serif CJK KR Medium:%
  language=KOR;script=hang;interhangul=-0.04em;interlatincjk=.125em;%
  charraise=-0.06em;+compresspunctuations" at 13.96pt
\def\titlefont{\latintitlefont\hangultitlefont}

\let\latinttitlefont=\ttitlefont
\sethangulfont\hangulttitlefont="Noto Sans Mono CJK KR:%
  language=KOR;script=hang;interhangul=-0.04em;interlatincjk=.125em;%
  charraise=-0.06em;+compresspunctuations" at 13.39pt
\def\ttitlefont{\latinttitlefont\hangulttitlefont}

% cwebmac sets strings, verbatim text and @d names in cmtex10, which has no
% Hangul; bring the monospace hangul (\tentz) into scope so that Korean inside a
% C string, inside \.{...} in a comment, or inside an @d name is typewriter as
% well, instead of falling back to the serif body face.  The original \. then
% sets cmtex10 for the Latin part as before.
% \leavevmode has to stay outside the group.  cwebmac's \. begins with it, so a
% paragraph that starts with \.{...} starts before any group is opened; if the
% group came first, \everypar would fire inside it and a hook that clears
% itself -- \everypar{\hskip-\parindent\everypar{}}, the way cwebman's \section
% macro unindents the paragraph after a heading -- would be restored when the
% group closed, and would go on unindenting later paragraphs.
\let\ttdot=\.
\def\.{\leavevmode\kottdot}
\def\kottdot#1{{\tentz\ttdot{#1}}}

% Running heads.  The group (chapter) title is upper-cased Latin, so Korean
% beside it wants the sans face (\ninegt, Noto Sans) rather than the serif body
% font; the document title is set as the author wrote it, in the same face.
\def\headfont{\eightrmfont\ninegt}
\def\lheader{\headertrue\mainfont\the\pageno\headfont\qquad\grouptitle
  \hfill\title\qquad\mainfont\topsecno} % top line on left-hand pages
\def\rheader{\headertrue\mainfont\topsecno\headfont\qquad\title\hfill
  \grouptitle\qquad\mainfont\the\pageno} % top line on right-hand pages

% Korean prose wants more leading than plain TeX's 12pt, and \tenpoint set it to
% 14pt above.  Displayed C code is Latin, though, and 14pt spreads it out; give
% the code parts (\B) a slightly tighter leading of their own.
\newskip\codebaselineskip \codebaselineskip=12.5pt
\let\Bnormal=\B
\def\B{\Bnormal\baselineskip=\codebaselineskip\relax}

% ---- repairing the discretionary breaks cweave puts inside strings ----------
% Every twenty bytes of a C string cweave inserts "}\)\.{", a discretionary
% backslash, counting bytes and not characters: a Korean character that straddles
% the boundary comes out cut in half, and the two halves are no longer valid
% UTF-8.  (Depending on where cweave then breaks its output line, the marker
% itself may be split as "}%" + "\)\.{" or as "}\)%" + "\.{".)  Since the damage
% is at the byte level, the repair has to happen before TeX decodes the line: a
% process_input_buffer callback removes the marker -- which rejoins the two
% halves -- but only when it really does sit inside a character, so that ordinary
% long strings keep their break points.
\begingroup
\catcode`\%=12 \catcode`\#=12 \catcode`\_=12 \catcode`\&=12
\catcode`\^=12 \catcode`\~=12 \catcode`\$=12
\directlua{
  local bs = string.char(92)
  local lb, rb = string.char(123), string.char(125)
  local mark = rb .. bs .. ")" .. bs .. "." .. lb
  local markpat = bs .. "%)" .. bs .. "%." .. lb
  local pending = nil
  local function contbyte(b) return b and b >= 128 and b < 192 end
  local function partial(s)
    for k = 1, 3 do
      local b = string.byte(s, -k)
      if not b or b < 128 then return nil end
      if b >= 192 then
        local n = 2
        if b >= 240 then n = 4 elseif b >= 224 then n = 3 end
        if n > k then return -k end
        return nil
      end
    end
    return nil
  end
  local function join(buf)
    local out, i = {}, 1
    while true do
      local a, b = string.find(buf, mark, i, true)
      if not a then break end
      if a > 1 and string.byte(buf, a - 1) >= 128
         and contbyte(string.byte(buf, b + 1)) then
        out[#out + 1] = string.sub(buf, i, a - 1)
      else
        out[#out + 1] = string.sub(buf, i, b)
      end
      i = b + 1
    end
    out[#out + 1] = string.sub(buf, i)
    return table.concat(out)
  end
  function cweb_repair_utf8(buf)
    if pending then
      local rest = string.match(buf, "^" .. markpat .. "(.*)$")
        or string.match(buf, "^" .. bs .. "%." .. lb .. "(.*)$")
      if rest and contbyte(string.byte(rest, 1)) then
        buf = pending .. rest
      else
        buf = pending .. rb .. buf
      end
      pending = nil
    end
    buf = join(buf)
    local body = string.match(buf, "^(.*)" .. rb .. bs .. "%)%%$")
      or string.match(buf, "^(.*)" .. rb .. "%%$")
    if body then
      local cut = partial(body)
      if cut then
        pending = string.sub(body, cut)
        buf = string.sub(body, 1, cut - 1) .. "%"
      end
    end
    return buf
  end
}
\endgroup
\directlua{
  if luatexbase then
    luatexbase.add_to_callback("process_input_buffer", cweb_repair_utf8,
      "cweb utf8 string repair")
  else
    callback.register("process_input_buffer", cweb_repair_utf8)
  end
}

% ---- Korean wording (cwebman's list of macros to redefine, point 14) --------
% The cross-reference notes printed under a code definition.  cwebmac's \note
% takes the wording first and the section numbers second; Korean puts the
% numbers first and the verb last, so \konote takes the wording in two pieces
% and sets  <before> <numbers> <after>.  Korean has no plural inflection, so the
% "section"/"sections" pairs read the same.
\def\konote#1#2#3.{\Y\noindent{\hangindent2em
    \baselineskip10pt\smallfont#1\ifacro{\pdfnote#3.}\else#3\fi#2.\par}}
\def\U{\konote{이 코드는 }{절에서 쓰인다}}
\def\Us{\konote{이 코드는 }{절에서 쓰인다}}
\def\A{\konote{이 코드는 }{절에도 정의되어 있다}}
\def\As{\konote{이 코드는 }{절에도 정의되어 있다}}
\def\Q{\konote{이 코드는 }{절에서 인용된다}}
\def\Qs{\konote{이 코드는 }{절에서 인용된다}}
\def\ch{\konote{변경 파일이 바꾼 절: }{}\let\*=\relax}
% Between the last two of several section numbers.  Korean would need 와 or 과
% depending on how the preceding numeral is pronounced, so a comma it is.
\def\ET{,~}
\def\ETs{,~}
% The name CWEAVE gives to the section that holds the preprocessor definitions.
\def\ATH{{\acrofalse\X\kern-.5em:전처리기 정의\X}}
% What @l prints in limbo.
\def\postATL#1 #2 {\bf 글자 \\{\uppercase{\char"#1}}는
   \tentex "#2"\bf 로 엮인다\egroup\par}
% The date, and the timestamp \datethis puts before the first section.
\def\today{{\number\year}년 {\number\month}월 {\number\day}일}
\def\datefont{\smallfont}
\def\dateat{\ }
\def\datethis{\def\startsection{\leftline{\datefont\today\dateat\hours}\bigskip
  \let\startsection=\stsec\stsec}}
\def\datecontentspage{\def\botofcontents{\vfill
   \centerline{\covernote}
   \bigskip
   \leftline{\datefont\today\dateat\hours}}}
% Headings: the list of section names, the table of contents, and its columns.
\def\namesheading{절 이름 목록}
\def\outsecname{절 이름 목록} % title of the section-names PDF bookmark
\def\contentsheading{목차}
\def\sectionword{절}
\def\pageword{쪽}
% The running head until the first starred section.  cwebmac has already put its
% own \gtitle into the first \mark, so the mark has to be set again.
\gtitle={\.{CWEB} 출력}
\mark{\noexpand\nullsec0{\the\gtitle}}

% ---- Korean PDF bookmarks ---------------------------------------------------
% A PDF outline title must be a UTF-16BE string, so Korean titles are converted
% by a small Lua helper, cweb_pdfutf16(s), which prints s as the body of a PDF
% literal string: the byte-order mark followed by each code unit as a pair of
% octal-escaped bytes (surrogate-paired above the BMP).  cwebmac's sanitizer has
% already put a backslash in front of the characters that are special inside a
% plain PDF string; octal escapes make that unnecessary, so the helper drops it.
%
% Two cautions about Lua inside \directlua: (1) the TeX-special characters
% % # _ & ^ ~ $ are made catcode 12 below (% is Lua's modulo operator, not a TeX
% comment; # is its length operator), and the backslash is produced as
% string.char(92) so that none appears in the source; and (2) \directlua turns
% the source's newlines into spaces, so a Lua "--" comment would swallow the
% rest of the chunk -- hence there are none here.
\begingroup
\catcode`\%=12 \catcode`\#=12 \catcode`\_=12 \catcode`\&=12
\catcode`\^=12 \catcode`\~=12 \catcode`\$=12
\directlua{
  local function oct(v)
    local d1 = math.floor(v / 64); v = v % 64
    local d2 = math.floor(v / 8); local d3 = v % 8
    return tostring(d1) .. tostring(d2) .. tostring(d3)
  end
  function cweb_pdfutf16(s)
    local bs = string.char(92)
    s = string.gsub(s, bs .. "(%p)", "%1")
    local t = { bs .. oct(254) .. bs .. oct(255) }
    for _, c in utf8.codes(s) do
      if c >= 65536 then
        c = c - 65536
        local w1 = 55296 + math.floor(c / 1024)
        local w2 = 56320 + (c % 1024)
        t[#t + 1] = bs .. oct(math.floor(w1 / 256)) .. bs .. oct(w1 % 256)
        t[#t + 1] = bs .. oct(math.floor(w2 / 256)) .. bs .. oct(w2 % 256)
      else
        t[#t + 1] = bs .. oct(math.floor(c / 256)) .. bs .. oct(c % 256)
      end
    end
    tex.sprint(-2, table.concat(t))
  end
}
\endgroup
% The sanitized title reaches us with \( \) \[ \] \/ standing for the characters
% that pdfTeX would have had to escape; make them ordinary again for Lua.  This
% has to happen before the title is handed to \pdfoutline, whose argument is
% merely expanded -- a \def inside it would be copied, not obeyed -- so
% \setbmtitle expands the title into \bmtitle first and \bmk does the rest.
\def\setbmtitle#1{\begingroup
  \def\({(}\def\){)}\def\[{[}\def\]{]}\def\/{/}%
  \xdef\bmtitle{#1}\endgroup}
\def\bmk{\directlua{cweb_pdfutf16("\luaescapestring{\bmtitle}")}}
% The bookmark for each starred section, written from the .toc file by \fin.
\def\writebookmarkline#1#2#3#4#5{\setbmtitle{#5}%
  \pdfoutline goto num #3 count -\expnumber{chunk#2.#3} {\bmk}}
% Open the bookmark (outline) pane when the document is first shown.
\def\startpdf{\pdfcatalog{/PageMode /UseOutlines}}

% ---- \fin and \con, localized ----------------------------------------------
% These two are cwebmac's own, with the English wording replaced and the
% bookmark titles passed through \setbmtitle.  cwebman (Appendix C, point 14)
% names them among the macros a translation has to redefine.
\def\fin{\par\vfill\eject % this is done when we are ending the index
  \ifpagesaved\null\vfill\eject\fi % output a null index column
  \if L\lr\else\null\vfill\eject\fi % finish the current page
  \ifpdftex \makebookmarks \fi
  \parfillskip 0pt plus 1fil
  \def\grouptitle{\namesheading}
  \let\topsecno=\nullsec
  \message{Section names:}
  \output={\normaloutput\page\lheader\rheader}
  \setpage
  \def\note##1##2##3.{\quad{\smallfont##1\ifacro{\pdfnote##3.}%
    \else{##3}\fi##2.}}
  \def\Q{\note{}{절에서 인용됨}} % crossref for mention of a section
  \def\Qs{\note{}{절에서 인용됨}}
  \def\U{\note{}{절에서 쓰임}} % crossref for use of a section
  \def\Us{\note{}{절에서 쓰임}}
  \def\I{\par\hangindent 2em}\let\*=*
  \ifacro \let\Xpdf\X
  \ifpdftex \pdfdest name {NOS} fith
    \setbmtitle{\outsecname}%
    \pdfoutline goto name {NOS} count -\secno {\bmk}
    \def\X##1:##2\X{\Xpdf##1:##2\X \firstsecno##1.%
      {\toksF={}\makeoutlinetoks##2\outlinedone\outlinedone}%
      \setbmtitle{\the\toksE}%
      \pdfoutline goto num \the\toksA {\bmk}}
  \fi\fi
  \readsections}

\def\con{\par\vfill\eject % finish the section names
  \rightskip 0pt \hyphenpenalty 50 \tolerance 200
  \setpage \output={\normaloutput\page\lheader\rheader}
  \titletrue % prepare to output the table of contents
  \pageno=\contentspagenumber
  \def\grouptitle{\contentsheading}
  \message{Table of contents:}
  \topofcontents \startpdf
  \line{\hfil\sectionword\hbox to3em{\hss\pageword}}
  \let\ZZ=\contentsline
  \readcontents\relax % read the contents info
  \botofcontents \end} % print the contents page(s) and terminate
